\documentclass[11pt,a4paper]{article}
\usepackage{dsmorgan}

\title{The Bessel Function of Order Zero}
\subtitle{Existence, uniqueness and an integral formula}
\author{Drew M. Schaffner}

\begin{document}
\maketitle

\section{Statement}

\begin{definition}[Bessel function]\label{def:j0}
  \begin{equation}\label{eq:series}
    J_0(z) = \sum_{m=0}^{\infty} \frac{(-1)^m z^{2m}}{4^m (m!)^2}
    = 1 - \frac{z^2}{4} + \frac{z^4}{64} - \dotsb .
  \end{equation}
\end{definition}

\begin{theorem}[Bessel's equation of order zero]\label{thm:main}
  \begin{enumerate}[label=(\alph*)]
    \item $J_0$ is entire and satisfies
      \begin{equation}\label{eq:bessel}
        z f''(z) + f'(z) + z f(z) = 0 \qquad (z \in \C).
      \end{equation}
    \item If $f$ is analytic on a disc about $0$ and satisfies \eqref{eq:bessel}
      there, then $f = f(0)\, J_0$.
    \item For every $z \in \C$,
      $J_0(z) = \dfrac{1}{\pi} \displaystyle\int_0^\pi \cos(z \sin\theta) \dd\theta$.
    \item $\abs{J_0(x)} \leq 1$ for every real $x$.
  \end{enumerate}
\end{theorem}

\section{Setting}

We use the following facts about power series \cite[ch.~5]{saff-snider}. A power
series is analytic inside its disc of convergence and may be differentiated term by
term there. If two power series agree on a disc about $0$, they have the same
coefficients.

\section{Idea of the proof}

Substituting a general power series into \eqref{eq:bessel} gives a two-step
recursion $n^2 a_n = -a_{n-2}$, with $a_1 = 0$. That recursion has exactly one
solution for each choice of $a_0$, and $a_0 = 1$ gives the coefficients of $J_0$.

\section{Proof}

\begin{lemma}[The recursion]\label{lem:recursion}
  Let $f(z) = \sum_{n \geq 0} a_n z^n$ converge on a disc about $0$. Then $f$
  satisfies \eqref{eq:bessel} on that disc if and only if
  \[
    a_1 = 0 \qquad\text{and}\qquad n^2 a_n + a_{n-2} = 0 \quad (n \geq 2).
  \]
\end{lemma}

\begin{proof}
  Differentiating term by term and collecting the coefficient of $z^{n-1}$,
  \[
    z f'' = \sum_{n \geq 2} n(n-1) a_n z^{n-1}, \qquad
    f' = \sum_{n \geq 1} n a_n z^{n-1}, \qquad
    z f = \sum_{n \geq 2} a_{n-2} z^{n-1}.
  \]
  So $z f'' + f' + z f = a_1 + \sum_{n \geq 2} (n^2 a_n + a_{n-2}) z^{n-1}$. This
  vanishes identically if and only if every coefficient is zero.
\end{proof}

\begin{proof}[Proof of (a) and (b)]
  Solve the recursion of \cref{lem:recursion}. From $a_1 = 0$ and
  $a_{2m+1} = -a_{2m-1}/(2m+1)^2$, every odd coefficient is zero. For the even ones,
  $a_{2m} = -a_{2m-2} / (4m^2)$, so by induction
  \[
    a_{2m} = \frac{(-1)^m a_0}{4^m (m!)^2}.
  \]
  For (b): $f$ has a power series at $0$, which must satisfy the recursion, so
  $f = a_0 J_0$ with $a_0 = f(0)$.

  For (a): the coefficients of $J_0$ satisfy the recursion with $a_0 = 1$, so it
  remains to see that \eqref{eq:series} converges everywhere. For $z \neq 0$, the
  ratio of consecutive terms is $\abs{z}^2 / (4(m+1)^2) \to 0$, so by the ratio test
  the radius of convergence is infinite. Then $J_0$ is entire, and
  \cref{lem:recursion} applies on every disc.
\end{proof}

\begin{lemma}[Wallis]\label{lem:wallis}
  For $m \geq 0$,
  $\displaystyle\int_0^\pi \sin^{2m}\theta \dd\theta = \pi \, \frac{(2m)!}{4^m (m!)^2}$.
\end{lemma}

\begin{proof}
  Let $I_m$ be the integral, so $I_0 = \pi$. For $m \geq 1$, integrate by parts with
  $u = \sin^{2m-1}\theta$ and $\dd v = \sin\theta \dd\theta$. The boundary terms
  vanish, so
  \[
    I_m = (2m - 1) \int_0^\pi \sin^{2m-2}\theta \cos^2\theta \dd\theta
        = (2m - 1)(I_{m-1} - I_m),
  \]
  that is, $I_m = \frac{2m-1}{2m} I_{m-1}$. The right-hand side of the lemma
  satisfies the same recursion, because
  \[
    \frac{(2m)!}{4^m (m!)^2} \bigg/ \frac{(2m-2)!}{4^{m-1} ((m-1)!)^2}
    = \frac{(2m)(2m-1)}{4m^2} = \frac{2m-1}{2m},
  \]
  and both sides equal $\pi$ when $m = 0$.
\end{proof}

\begin{proof}[Proof of (c)]
  Fix $z$. The series $\cos(z\sin\theta) = \sum_m (-1)^m z^{2m} \sin^{2m}\theta / (2m)!$
  converges uniformly in $\theta \in [0, \pi]$, since its $m$th term is at most
  $\abs{z}^{2m} / (2m)!$. So we may integrate term by term and apply \cref{lem:wallis}:
  \[
    \frac1\pi \int_0^\pi \cos(z \sin\theta) \dd\theta
    = \sum_{m=0}^{\infty} \frac{(-1)^m z^{2m}}{(2m)!} \cdot \frac{(2m)!}{4^m (m!)^2}
    = J_0(z). \qedhere
  \]
\end{proof}

\begin{proof}[Proof of (d)]
  For real $x$ the integrand in (c) is a real number in $[-1, 1]$, so
  $\abs{J_0(x)} \leq \frac1\pi \int_0^\pi 1 \dd\theta = 1$.
\end{proof}

\section{Remarks}

\begin{remark}[An energy proof of (d)]
  Part (d) also follows from the equation alone. For real $x > 0$, put
  $E = J_0^2 + (J_0')^2$. By \eqref{eq:bessel},
  $J_0'' = -J_0 - J_0'/x$, so
  \[
    E' = 2 J_0' (J_0 + J_0'') = -\frac{2 (J_0')^2}{x} \leq 0 .
  \]
  Hence $E(x) \leq E(0) = 1$, so $\abs{J_0(x)} \leq 1$. Since $J_0$ is even, the same
  holds for $x < 0$. Physically, the term $f'/x$ acts as friction, so the
  oscillation loses energy and its amplitude decays.
\end{remark}

\begin{remark}
  Away from $0$, \eqref{eq:bessel} is a second-order linear equation, so it has a
  two-dimensional space of solutions on any disc not containing $0$. By (b), only the
  multiples of $J_0$ extend analytically to $0$. Every independent second solution is
  singular at the origin; the standard choice is the Bessel function $Y_0$ of the
  second kind, which has a logarithmic singularity there.
\end{remark}

% Verified: archive.org record fundamentalsofco0000saff_c0g0. Cite by chapter only.
\begin{thebibliography}{9}
  \bibitem{saff-snider} E.~B. Saff and A.~D. Snider,
    \emph{Fundamentals of Complex Analysis with Applications to Engineering
    and Science}, 3rd ed., Prentice Hall, Upper Saddle River, NJ, 2003.
\end{thebibliography}

\end{document}
